\documentclass[../../../uem_phy_std_a.tex]{subfiles}

\begin{document}
    \begin{enumerate}
        \item AB間を流れる電流の大きさを$I$とする．キルヒホッフ則より，
        
            \myspaceb%
            $\displaystyle%
                3RI + 2RI = V_0 \qquad \therefore \; I = \myfracc{V_0}{5R} \;.
            $\\
            よって，
            
            \myspaceb%
            $\displaystyle%
                V_{2} = 3RI = \uwave{\myfraca{3}{5}V_0} \;.
            $
        \item R$_1$の電流を$I$（右向き），R$_2$の電流を$i$（右向き）とすると，
            R$_3$の電流は$I - i$（右向き）である．キルヒホッフ則より，
        
            \myspaceb%
            $\displaystyle%
                \Biggl\{%
                \begin{myaligna}
                    & Ri = 2R(I - i) \;, \\
                    & Ri + 2RI = V_0
                \end{myaligna}
                \qquad \therefore \; i=\uwave{\myfracc{V_0}{4R}} \;,
                \quad I = \myfracc{3V_0}{8R} \;.
            $
        \item R$_2$の電流を$i$（右向き）とすると，R$_3$の電流は$I-i$（右向き）である．
            キルヒホッフ則より，

            \myspaceb%
            $\displaystyle%
                \Biggl\{%
                \begin{myaligna}
                    & Ri + RI = 2V_0 \;, \\
                    & Ri = 2R(I - i) + V_0
                \end{myaligna}
                \qquad \therefore \; i=\myfracc{V_0}{R} \;,
                \quad I = \uwave{\myfracc{V_0}{R}} \;.
            $
        \item R$_1$の電流を$i$（右向き），R$_5$の電流を$j$（上向き）とすると，
            R$_2$の電流は$i+j$（右向き），R$_3$の電流は$I-i$（右向き），
            R$_4$の電流は$I-i-j$（右向き）である．キルヒホッフ則より，
                    
            \myspaceb%
            $\displaystyle%
                \Biggl\{%
                \begin{myaligna}
                    & R(i+j) + 2Ri = R(I-i-j) + R(I-i) = V_0 \;,\\
                    & 2Ri = Rj + R(I-i) \;.
                \end{myaligna}
            $\\
            よって，
            
            \myspaceb%
            $\displaystyle%
                i = \myfracc{4V_0}{13R} \;, \quad j = \myfracc{V_0}{13R} \;,
                \quad I = \uwave{\myfracc{11V_0}{13R}} \;.
            $
    \end{enumerate}
\end{document}
